\chapter{总结与展望}

\section{工作总结}

本文围绕处理器微架构侧信道漏洞安全测试，使用Flush+Reload分析AES加密算法在多密钥长度下的攻击可行性和缓解措施的有效性。

在攻击框架上，本实验构建了一套支持AES-128/192/256三种密钥长度的统一测试框架，采用模块化设计。整个流程划分为六个功能模块：共享内存管理、Victim进程模拟、Attacker进程探测、密钥恢复、量化评估、自动化实验。Victim和Attacker采用父子进程模型，管道同步通信确保攻击时序可控。共享内存模块使用POSIX shm\_open和mmap，T表数据映射到双方地址空间，为Flush+Reload探测提供共享页面。自动化实验支持批量配置和结果可视化，不同参数组合的对比实验可高效完成。

密钥恢复策略为三种AES密钥长度分别做了适配。AES-128通过排除法恢复最后一轮子密钥$K_{10}$，再通过密钥扩展逆过程推导原始密钥$K_0$。实验中5000个样本即可稳定恢复完整128位密钥。AES-192和AES-256需采用两阶段：阶段一恢复最后一轮子密钥（AES-192恢复$K_{12}$，AES-256恢复$K_{14}$）。阶段二复用同一批Flush+Reload监控数据，结合等效解密中间值$IV_1$，对等效解密轮密钥$Kd_1$的每个字节直接执行排除法，恢复$Kd_1$后经MixColumns变换得到加密倒数第二轮子密钥。复用监控数据的优势在于无需重复采样——AES-128、AES-192、AES-256分别使用5000、9000、15000个样本。

评估体系涵盖基础指标、时间统计指标、泄露带宽和攻击成功率。基础指标包括样本总数、命中次数、未命中率。时间统计指标采用分位数（中位数、Q1、Q3、P5、P95）替代均值和标准差——分位数对用户态高精度时间测量中常见的异常值具有更强的鲁棒性。泄露带宽定义为$BW = 4 \cdot MI_{binary} \cdot f_{sample}$，其中$MI_{binary}$为单次缓存行观测的互信息，系数4对应每次采样对4个缓存行各进行一次独立观测，以互信息与采样频率的乘积从信息泄露速率角度刻画侧信道威胁。

缓解措施评估中，对噪声注入和缓存刷新两种方法进行了系统化评估。基准测试中三种AES变体在无缓解措施的情况下成功率均为100\%。缓存刷新使成功率与泄露带宽均降至零，防护彻底。噪声注入方面发现：密钥越长，加密轮数越多，$p_{miss}$越低，每阶段有效信息量越少，加之两阶段攻击的32字节乘积级联放大效应，对噪声亦越敏感。

\section{未来工作展望}

跨平台验证。框架目前在x86上已验证通过，但ARM和RISC-V架构的缓存层次结构存在差异，T表的缓存行布局可能需要重新适配。虚拟化环境和容器环境中，Hypervisor或容器运行时可能对CLFLUSH指令和共享内存施加额外限制，可行性亦需重新评估。

缓解措施的优化。噪声注入和缓存刷新各有优劣——缓存刷新防护彻底但性能开销大，噪声注入开销可控但防护不彻底。自适应缓解是一个值得探索的方向：根据运行时检测到的侧信道威胁指标动态调整噪声强度，在安全与性能之间寻求动态平衡。
